\documentclass{article}
\usepackage[legalpaper, portrait, margin=0.5in]{geometry}
\usepackage{amsmath, cancel, amssymb}

\title{CSC226 Exercise 2}
\author{Dryden Bryson}
\date{\today}

\begin{document}

\maketitle
\newpage

\section*{Exercise 2.}

The overall running time of Quickselect with BFPRT partitioning with subarrays of size 7 can be modelled by the following recurrence:

$$T(n)=\begin{cases}
c_{b} & \text{if} \;n \leq 1\\
T(\lfloor \frac{5n}{7} \rfloor ) + T(\lfloor \frac{n}{7} \rfloor ) + c_{1}n+c_{2} & \text{if} n \geq 2
\end{cases}$$

Where:
\begin{itemize}
    \item $T(\lfloor \frac{5n}{7} \rfloor )$ is the worst case recursive quick sort call, this is found since there is $\lfloor \frac{n}{7} \rfloor $ sub-arrays when selecting our pivot. And once we find our median of medians, half of the medians ($\lfloor \frac{n}{14} \rfloor $) are guaranteed to be bigger, and since each median has 4 elements which are $\leq$ to the median, when we run recurse upon quicksort we are guaranteed to be getting rid of $\lfloor 4\cdot \frac{n}{14} \rfloor $ Thus we recurse on, in the worst case: $$\lfloor n-\frac{4n}{14} \rfloor = \lfloor n-\frac{2n}{7} \rfloor =\lfloor \frac{5n}{7} \rfloor $$
    of the array.
    \item $T(\lfloor \frac{n}{7} \rfloor )$ is for finding the medians of our $\lfloor \frac{n}{7} \rfloor $ sub-arrays.
    \item $c_{1}n$ linear constant work, i.e sorting sub arrays, partitioning array.
    \item $c_{2}$ constant work
\end{itemize}
Let $c=c_{b}+c_{1}+c_{2}$ We form $T'(n)$ such that $T'(n) \geq T(n)$ for simplification, we have:

$$
T'(n)=
\begin{cases}
c & \text{if}\;n\leq 1\\
T(\lfloor \frac{5n}{7} \rfloor )+ T(\lfloor \frac{n}{7} \rfloor) + cn & \text{if}\;n \geq 2
\end{cases}
$$
\\
We will proceed with the proof by applying induction on $n$:\\
\textbf{Claim:} $T'(n) \leq 7cn$ for $n\geq 1$\\\\
\textbf{Proof:} By induction on $n$\\\\
\textbf{Basis:} When $n=1$, $T'(n)=c=cn\leq 7cn$\\\\
\textbf{Induction Hypothesis:} Suppose $T'(n)\leq 7cn$ for $n \in \left\{ 1,2,\dots,k \right\} $ for some $k\geq 1$\\\\
\textbf{Induction Step:} Consider $T'(n+1)$
\begin{table}[htp]
\centering
\begin{tabular}{ccll}
$T'(n+1)$   & $=$  & $T'(\lfloor \frac{5(n+1)}{7} \rfloor )+T'(\lfloor \frac{n+1}{7} \rfloor )+c(n+1)$  &   \\
   &  $\leq$ & $7c(\lfloor \frac{5(n+1)}{7} \rfloor )+7c(\lfloor \frac{n+1}{7} \rfloor )+c(n+1)$  & \text{by the induction hypothesis}  \\
   &  $\leq$ &  $7c( \frac{5(n+1)}{7}  )+7c( \frac{n+1}{7}  )+c(n+1)$  &  \text{Since} $\forall x \in \mathbb{R}, \lfloor x \rfloor  \leq x$  \\
   & $=$  &  $ \frac{\cancel{7}c\cdot 5(n+1)}{\cancel{7}}  )+\frac{\cancel{7}c \cdot( n+1)}{\cancel{7}}  +c(n+1)$  &   \\
   & $=$  &  $5c(n+1)+c(n+1)+c(n+1)$ &   \\
   & $=$  &  $7c(n+1)$ &   \\
   &   &   &   \\
\end{tabular}
\end{table}\\

Since we have shown that $T'(n+1)\leq 7c(n+1)$ our claim holds and thus, by induction, our claim also holds for all values of $n$. We can then conclude that $T'(n)=\Theta(n)$ and consequently, $T(n)=\Theta(n)$.


\end{document}